% ============================================================
%  INTISARI / ABSTRACT (English)
% ============================================================
 
Most websites still run PHP applications on version 7.x or older that no longer receive official security updates, leaving systems exposed to an ever-growing accumulation of vulnerabilities with no possibility of remediation. Manual migration across hundreds of code files is impractical, and current transformation tools do not include security checks as an integrated part of the migration workflow. Organizations operating under the ISO/IEC 27001:2022 standard must also perform security audits as a separate step after migration is complete, adding to the operational burden.
 
\hspace{1cm} This research builds a Python-based platform that integrates automatic conversion of legacy PHP web applications to a configurable target PHP version with ISO/IEC 27001:2022 secure coding compliance checks in a single, locally executable open-source workflow. The platform uses Rector for AST-based code transformation, Semgrep for security scanning, PHPStan for post-conversion validation, and five open-source LLMs running locally via Ollama---DeepSeek Coder 6.7b, Qwen2.5-Coder 7b, CodeLlama 7b, Mistral 7b, and Llama 3.1 8b---for semantic analysis and vulnerability remediation suggestions. Case studies were conducted on PHP codebases from FT UGM (Dashboard TCK, 245 files) and CBT DTI UGM (326 files).
 
\hspace{1cm} Rector achieved a conversion rate of 97.6\% on the FT UGM dataset and 91.7\% on the CBT DTI UGM dataset, with a post-conversion syntax validity rate of 100\% in both case studies. A value of $\Delta F = 0$ confirms that the conversion process introduced no new vulnerabilities. The automatically generated JSON reports include per-finding mappings to the relevant ISO/IEC 27001:2022 Annex A controls and can be used directly as audit documentation. From a comparative experiment across five LLM models, Qwen2.5-Coder 7b is recommended for integration into automated systems, as it consistently achieved a Syntax Validity Rate (SVR) of 100\% across all experimental conditions. This platform demonstrates that the integration of automated code conversion and ISO/IEC 27001:2022 compliance checking can be executed entirely locally without transmitting source code to any external service.
 
\noindent{Keywords}: PHP migration, secure coding, ISO/IEC 27001:2022, large language model, AST code transformation